% Shared preamble for the "Values in Ranking" reddit-post ranking handouts.
%
% Self-contained (does not \input the CS 377 worksheet template) so this
% directory can be copied or archived on its own. Mirrors the look of
% misc/worksheets/00-template.tex (course line, big title, member/section
% tables) plus a ranking table for the reddit posts.
%
% Built with generate_ranking_pdfs.py, which writes one .tex per version and
% compiles it with latexmk.

\documentclass[11pt]{article}

\usepackage[T1]{fontenc}
\usepackage[utf8]{inputenc}
\usepackage{lmodern}
\usepackage[margin=0.9in, top=0.6in, bottom=0.6in]{geometry}
\usepackage{microtype}
\usepackage{array}
\usepackage{tabularx}
\usepackage{longtable}
\usepackage[table]{xcolor}
\usepackage{calc}

\setlength{\parindent}{0pt}
\setlength{\parskip}{0.5\baselineskip}
\pagestyle{empty}
\raggedbottom

\renewcommand{\tabularxcolumn}[1]{m{#1}}

% ---------------------------------------------------------------------------
% Metadata
% ---------------------------------------------------------------------------
\newcommand{\wscourseval}{CS 377}
\newcommand{\wstitleval}{Worksheet}
\newcommand{\wscourse}[1]{\renewcommand{\wscourseval}{#1}}
\newcommand{\wstitle}[1]{\renewcommand{\wstitleval}{#1}}

\newcommand{\wssections}{\wssectiontime{12:30pm} & \wssectiontime{2:00pm}}
\newcommand{\wssectioncount}{2}

% ---------------------------------------------------------------------------
% Header pieces
% ---------------------------------------------------------------------------
\newcommand{\wscourseline}{%
  {\raggedleft\small\wscourseval\par}%
  \vspace{0.4\baselineskip}%
}

\newcommand{\wsbigtitle}{%
  {\fontsize{22}{26}\selectfont\bfseries\wstitleval\par}%
  \vspace{0.3\baselineskip}%
}

\newcommand{\wsmembertable}{%
  \begin{tabularx}{\textwidth}{|X|X|}
    \hline
    \rule{0pt}{2.4\baselineskip} & \\
    \hline
    \rule{0pt}{2.4\baselineskip} & \\
    \hline
  \end{tabularx}%
}

\newcommand{\wsvcell}[2]{\parbox[c][#1][c]{\linewidth}{#2}}

\newlength{\wssectionrowht}
\newcommand{\wssectiontime}[1]{\wsvcell{\wssectionrowht}{\centering #1}}
\newcommand{\wssectiontable}{%
  \setlength{\wssectionrowht}{2.6\baselineskip}%
  \begin{tabularx}{\textwidth}{|>{\bfseries\small}m{1.1in}|*{\wssectioncount}{X|}}
    \hline
    \wsvcell{\wssectionrowht}{\raggedleft Section\newline(circle one)} & \wssections \\
    \hline
  \end{tabularx}%
}

\newcommand{\wsheader}{%
  \wscourseline
  \wsbigtitle
  Group members:\par
  \vspace{0.2\baselineskip}
  \wsmembertable
  \vspace{0.6\baselineskip}
  \wssectiontable
  \vspace{0.4\baselineskip}
}

\newcounter{wsquestion}
\newcommand{\wsprompt}[1]{%
  \stepcounter{wsquestion}%
  \par\noindent(\thewsquestion)~#1\par
}

% ---------------------------------------------------------------------------
% Ranking-table pieces
% ---------------------------------------------------------------------------

% \rkinstructions{rank max} -- the two-line instructions under the title.
\newcommand{\rkinstructions}[1]{%
  Discuss the following posts with your table and propose a ranking (1-#1) for each post.\par
  Rank 1 means that the post is placed highest and should be seen by the most people.\par
  \vspace{0.4\baselineskip}
}

% \rktable{rows} -- rows is a sequence of "title & preview & tag & \\ \hline"
% lines, generated by generate_ranking_pdfs.py.
% \rowcolors's alternating fill is a document-wide toggle, not scoped to one
% table, so \rktable wraps it in a group -- otherwise the stripes would leak
% into the member/section tables on the reflection page.
\newcommand{\rktable}[1]{%
  \begingroup
  \setlength{\tabcolsep}{4pt}
  \rowcolors{2}{gray!8}{white}
  \begin{longtable}{|>{\raggedright\arraybackslash}m{1.5in}|>{\raggedright\arraybackslash}m{2.6in}|>{\raggedright\arraybackslash}m{0.8in}|>{\raggedright\arraybackslash}m{0.95in}|}
    \hline
    \textbf{Title} & \textbf{Post Preview} & \textbf{Tag} & \textbf{Proposed\newline Ranking\newline(1=Top)} \\
    \hline
    \endfirsthead
    \hline
    \textbf{Title} & \textbf{Post Preview} & \textbf{Tag} & \textbf{Proposed\newline Ranking\newline(1=Top)} \\
    \hline
    \endhead
    \hline
    \endfoot
    #1
  \end{longtable}%
  \endgroup
}
